\chapter{The Composition Failure Theorem}
\label{ch:composition-failure}

\begin{quotation}
\noindent\textit{Synopsis.} This chapter presents the monograph's central
new result. The Composition Failure Theorem demonstrates that a sequence of
locally rational transformations can collectively produce a globally
insufficient outcome, with no stage behaving negligently, incompetently, or
irrationally: failure emerges because each stage optimizes for its own
local task while future tasks remain unknown. This formalizes a pattern
familiar from archival systems, institutions, scientific communication, and
information-processing pipelines, in which a distinction that seemed
irrelevant at one stage becomes essential later.
\end{quotation}

\newresult

\section{Making Local Admissibility Precise}

Before the theorem can be stated, the phrase it turns on --- a transition
being ``locally admissible with respect to its own stage's task''
--- needs a definition sharper than \cref{def:budget-discard}'s general
schema alone provides. \cref{def:budget-discard} says a distinction is
dropped when funding it exceeds budget or when it interferes too
strongly with a higher-priority distinction; it does not say how
priority itself is supposed to track a stage's task. The following
definition closes that gap, and everything downstream --- the theorem,
its corollaries, and the proof of Appendix~\ref{app:proof} --- depends on
exactly this closure.

\begin{definition}[Locally Admissible Transition]
\label{def:locally-admissible}
A transition $B_i$ is locally admissible with respect to $\tau_i$ if its
discard set $D_i$ is produced by a construction satisfying
\cref{def:budget-discard} in which every distinction required by $\tau_i$
(\cref{def:benign-loss}) is ranked strictly above every distinction not
required by $\tau_i$.
\end{definition}

This definition is deliberately weak in one respect and strong in
another. It is weak in that it does not choose between
\cref{con:greedy-discard} and \cref{con:optimal-discard}, or any other
construction satisfying \cref{def:budget-discard}: any of them counts as
locally admissible provided the priority ordering they operate on is
correctly aligned with $\tau_i$. It is strong in that this alignment is
not optional --- a transition that funds a distinction $\tau_i$ does not
need ahead of one it does need fails to be locally admissible, regardless
of how sophisticated the underlying discard construction is. A transition
satisfying \cref{def:locally-admissible} cannot be accused of
irrationality relative to $\tau_i$: it funds everything $\tau_i$ needs, in
full, before spending anything on distinctions $\tau_i$ does not need.
This is exactly the standard the theorem needs every stage to meet, since
the entire point of the result is that failure can occur even when no
stage falls short of it.

\section{Statement}

\begin{theorem}[Composition Failure]
\label{thm:composition-failure}
Let $A_0 \xrightarrow{B_0} \cdots \xrightarrow{B_{n-1}} A_n$ be a budget
relay (\cref{def:budget-relay}) in which each transition $B_i$ is locally
admissible (\cref{def:locally-admissible}) with respect to its own
stage's task $\tau_i$. Then there exist open task spaces $T$
(\cref{prop:task-unavailability}) containing some $\tau^* \notin \{\tau_0,
\ldots, \tau_{n-1}\}$ such that $A_n$ is insufficient
(\cref{def:sufficiency}) for $\tau^*$, even though no individual
transition was locally irrational relative to its own $\tau_i$.
\end{theorem}

Each clause of this statement is load-bearing, and it is worth checking
what happens if any one is dropped. Without \cref{def:budget-relay}'s
independence condition, the chain reduces to a single observer's economy
and \cref{thm:redistribution} governs it instead, as
Appendix~\ref{app:comparison} shows concretely. Without local
admissibility, the theorem would be uninteresting: of course an
incompetently-run relay can end up insufficient. Without $T$ being open,
$\tau^*$ could in principle have been anticipated and funded for, and the
failure would trace back to ordinary negligence rather than to the
structure of uncoordinated composition itself. The theorem's content is
precisely that none of these outs are available: the relay is
uncoordinated by definition, every stage did what it was supposed to do,
and the task that breaks it was not available to be planned for.

\section{Proof Strategy}

The full proof is given in Appendix~\ref{app:proof}; this section
previews its shape so the construction does not appear unmotivated when
read there.

\begin{proofsketch}
The proof constructs an explicit witness relay using the five-distinction
instance introduced in \cref{sec:worked-relay} and reused in
\cref{sec:redistribution-worked}, together with a family of
\emph{recovery tasks} $\tau_d$ --- one for each named distinction $d$,
asking only whether $d$ holds in the source
(\cref{def:recovery-task}) --- which serves as the open task space $T$.
A stage whose required distinctions exactly, or near-exactly, exhaust its
local budget is shown (\cref{lem:exhaustion}) to drop every non-required
distinction \emph{regardless of which construction satisfying
\cref{def:budget-discard} is used}, which is what allows the witness to
be locally admissible in the strong sense of \cref{def:locally-admissible}
without leaving room for a more clever discard policy to have rescued the
dropped distinction. Some distinction $d^*$ -- concretely, $d_5$ in the
worked witness -- is required by neither stage's task and is dropped at
the earliest stage as a direct consequence. By
\cref{prop:monotonicity}, $d^*$ remains unavailable at every later stage
absent root recovery (\cref{prop:root-recovery}), which the witness does
not invoke. Taking $\tau^* = \tau_{d^*}$ then gives insufficiency
directly, by \cref{lem:recovery-sufficiency}: a rendering is sufficient
for a recovery task exactly when it has not discarded the distinction in
question, and $A_n$ has.
\end{proofsketch}

This preview differs in one respect from an earlier, less precise sketch
of the same argument, which described $d^*$ as dropped because it
``interferes with a distinction $\tau_i$ does require.'' The actual
mechanism proved in Appendix~\ref{app:proof} is sharper and does not rely
on interference at all: $d^*$ is dropped because, after funding every
required distinction in full, no budget capable of funding $d^*$ remains
--- a conclusion that holds regardless of interference structure and
regardless of which construction satisfying \cref{def:budget-discard}
produced the allocation. The interference-based mechanism remains
available in principle for other witnesses, but it is not what the
canonical witness of this monograph uses, and the sketch above is
corrected accordingly.

\section{Corollaries}

\begin{corollary}[Local Improvement Does Not Guarantee Global Sufficiency]
\label{cor:local-improvement}
Improving any single stage's local budget cannot, by itself, eliminate
all instances of the failure mode exhibited in
\cref{thm:composition-failure}.
\end{corollary}

\begin{proofsketch}
Verified directly against the witness in Appendix~\ref{app:proof}:
increasing $\beta_1$ cannot restore a distinction already absent from
$A_1$ by the time stage $1$ acts, since \cref{prop:monotonicity} forbids
recovery absent root access. Only an improvement at the specific stage
where the loss occurred could have prevented it, and even that requires
knowing, at that stage, that the eventually-relevant task would need the
distinction in question -- precisely the knowledge
\cref{prop:task-unavailability} shows is generally unavailable.
\end{proofsketch}

\begin{corollary}[No Free Lunch for Relays]
\label{cor:no-free-lunch}
A budget relay with $n$ uncoordinated stages has cumulative discard
(\cref{def:cumulative-discard}) that is non-decreasing in $n$
(\cref{prop:monotonicity}), realized with equality across arbitrarily
many additional pass-through stages
(\cref{prop:extension}); the coordination mechanisms of
Chapter~\ref{ch:coordination} can narrow, but not eliminate, the
resulting insufficiency (\cref{prop:no-coordination-universal}).
\end{corollary}

\begin{proofsketch}
The non-decrease is \cref{prop:monotonicity} directly, and
Appendix~\ref{app:proof} shows this bound is achieved with equality
rather than merely respected, since pass-through stages contribute no
further discard. This corollary should not be read as a claim that
cumulative discard \emph{risk} grows in some stronger, quantitative
sense across relays in general; that stronger claim is not established
here and is recorded as open in Appendix~\ref{app:open-questions}. What
is established is the existence half only: discard does not shrink, and
can accumulate to the full extent this chapter's witness demonstrates.
\end{proofsketch}

\section{Relation to the Redistribution Theorem}

\cref{thm:composition-failure} is not a special case of
\cref{thm:redistribution}, and is not derivable from it without the
budget-relay machinery of Chapter~\ref{ch:budget-relay}. The
Redistribution Theorem concerns simultaneous competition for one
observer's budget; the Composition Failure Theorem concerns a single
distinction's fate across a sequence of independent, non-communicating
budgets. Appendix~\ref{app:comparison} develops this contrast concretely
rather than only structurally: pooling the witness's two stage budgets
into a single shared total is shown there
(\cref{prop:pooling-restores}) not merely to make
\cref{thm:redistribution} applicable again, but to actually succeed at
funding the distinction the separated relay was forced to drop, using
the identical total nominal resource. The insufficiency this chapter's
theorem exhibits is therefore a consequence specifically of
\cref{def:budget-relay}'s prohibition on transferring unspent capacity
across stages, not of any absolute scarcity of resource -- confirming
that the two theorems occupy genuinely disjoint settings rather than one
subsuming the other.
